\documentclass[10pt,a4paper]{amsart}
\usepackage[margin=19mm]{geometry}
\usepackage{amsmath,amssymb,mathtools}
\usepackage[scr=rsfs,cal=euler]{mathalfa}
\usepackage{xfrac}
\usepackage[unicode,hidelinks,bookmarks=false]{hyperref}
\newcommand{\ie}{i.e.\ }
\newcommand{\cf}{cf.\ }
\newcommand{\resp}{respectively\ }
\newcommand{\Subsection}{\subsection}
\providecommand{\nobreakdash}{\nobreak\hbox{-}}
\setlength{\emergencystretch}{2em}
\multlinegap0pt
\usepackage{amssymb}
\usepackage{xcolor}
\newtheorem{thm}{Theorem}
\newcommand{\C}{\mathbb C}
\newcommand{\Dotprod}[2]{\big\langle\!\:\!\!\big\langle #1, #2 \big\rangle\!\:\!\!\big\rangle}
\newcommand{\Int}{\mathbb Z}
\newcommand{\abs}[1]{\vert #1 \vert}
\newcommand{\bcode}[1]{\texttt{#1}}
\newcommand{\bs}[1]{{\boldsymbol{#1}}}
\newcommand{\conv}{\ast}
\newcommand{\dimin}[2]{{ #1 }^{\gr{#2}}}
\newcommand{\dotprod}[2]{\langle\!\langle #1, #2 \rangle\!\rangle}
\newcommand{\gr}[1]{( #1 )}
\newcommand{\hatbsf}{{\kern+0.25em}\hat{{\kern-0.25em}\bs{f}}}
\newcommand{\ip}[2]{\langle #1, #2 \rangle}
\newcommand{\klth}{$(k,l)$th}
\newcommand{\mth}{$m$th}
\newcommand{\qth}{$q$th}
\newcommand{\vect}[1]{[ #1 ]}
\title{Statistics of a multi-factor function\\ from its Fourier transform}
\author{Matthew A. Herman, Stephen Doro}
\date{Source: arXiv:2605.02248v1 (CC BY 4.0)}
\begin{document}
\maketitle

\noindent\emph{Source and licence.} Statistics of a multi-factor function\\ from its Fourier transform. Version arXiv:2605.02248v1 (CC BY 4.0), distributed by arXiv under the Creative Commons Attribution 4.0 International licence, http://creativecommons.org/licenses/by/4.0/, as recorded in the arXiv licence metadata for this exact version and shown on the abstract page https://arxiv.org/abs/2605.02248v1. Full licence notice: \texttt{licenses/arxiv-2605.02248-CC-BY-4.0.txt}.\par\smallskip

\noindent\textit{Editorial scope.} Counted unit: the article's thm Thm:mthMoment\_a (source0.tex lines 607--617). The article's complete original proof of it is delivered with it (source0.tex lines 619--633, one \texttt{proof} environment of 15 lines). No mathematics is written, repaired or re-derived: the statement and the proof are the article's own lines, in the article's own notation, and no retained line is substituted or re-flowed.  Retained uncounted as the source-local context the proof needs: lines 187--189; lines 224--226; lines 248--250; lines 319--321; lines 341--344; lines 354--356; lines 372--375; lines 381--389; lines 425--427; lines 527--529; lines 541--556.  Nothing was omitted from any retained span, so the package carries no omission record.  No retained line contains a citation, so the wrapper carries no bibliography and no bibliography entry.  No retained line contains a figure environment, an image-inclusion command or drawing code, so no image is delivered and the package declares no asset dependency.  Editorial formatting only, in addition: an AMS \texttt{amsart} preamble that restates the article's own declarations and macros used by the retained lines, namely the environments \texttt{thm} on separate counters, together with the standard packages those lines need.  The retained environments print in this document as Theorem 1 in order of appearance, whereas the article prints the same items with the numbering of its own full document.  The complete native source of the article as distributed in the arXiv e-print for this exact version is bundled unchanged as \texttt{sources/199708/source0.tex}.
\begin{equation} \label{Def:mthMoment_a}
\dimin{\mu}{a}_m\gr{\bs{f}} \,:=\; \frac{1}{\abs{G}} \sum_{i\in G} \gr{f_i-a}^m.
\end{equation}

\begin{equation} \label{Eq:FiniteAbelianGroup_G}
G \,=\, \Int_{N_1} \times \Int_{N_2} \times \cdots \times \Int_{N_n}
\end{equation}

\begin{equation} \label{Def:dotProduct}
\qquad\dotprod{\bs{f}}{\bs{g}} \,:=\; \sum_{i\in G} f_i \:\! g_i \,=\, \ip{\bs{f}}{\overline{\bs{g}}}, \qquad \text{for any } \bs{f},\bs{g}\in L\gr{G}.
\end{equation}

\begin{equation} \label{Def:F-pair}
\bs{f} \;\stackrel{\mathcal{F}}\rightleftharpoons\; \hatbsf
\end{equation}

\begin{equation} \label{Eq:f^m=*^mfhat}
\bs{f}^m \;\stackrel{\mathcal{F}}\rightleftharpoons\;
\conv^m \hatbsf
\end{equation}

\begin{equation} \label{Eq:Parseval_fg_dotprod}
\dotprod{\bs{f}}{\bs{g}} \,=\ \abs{G} \:\!\dotprod{\hatbsf}{\hat{\bs{g}}_\ominus}
\end{equation}

\begin{equation} \label{Def:a-diminished_F-pair}
\bs{f}-a\bs{1} \;\;\stackrel{\mathcal{F}}\rightleftharpoons\;\;
\dimin{\hatbsf}{a}
\end{equation}

    \begin{equation} \label{Def:a-diminished_Xform}
    \dimin{\hatbsf}{a} :=\, \hatbsf - a\bs{e}_0 \,=\,
    \left\{
      \begin{array}{ll}
        \dimin{\hat{f}}{a}_0 = \mu-a, & \quad\hbox{if $j=0$;} \\[4pt]
        \dimin{\hat{f}}{a}_j = \hat{f}_j,    & \quad\hbox{otherwise.}
      \end{array}
    \right.\;
    \end{equation}

\begin{equation} \label{Eq:C_fC_g_ij}
\bs{C}_{\:\!\!\bs{f}}\:\!\bs{C}_{\;\!\!\bs{g}}\gr{i,j} \,=\, \dotprod{\mathcal{S}_i\bs{f}_{\!\ominus}}{\mathcal{S}_j\bs{g}} \,=\, \sum_{l\in G} f_{i\ominus l} \:\! g_{l\ominus j} \,=\, \sum_{\substack{p,q\in G\\[2pt] p\oplus q = i \ominus j}} f_p \;\! g_q
\end{equation}

\begin{equation} \label{Eq:C_*^mf=C_f^m}
\bs{C}_{\gr{\:\!\!\conv^m\:\!\!\bs{f}}} \,=\, \bs{C}_{\:\!\!\bs{f}}^m
\end{equation}

\subsection{The \mth\ moment from Fourier coefficients}
Recall the \mth\ general moment of~$\bs{f}$ from~\eqref{Def:mthMoment_a} and assume functions $\bs{f} \! \stackrel{\mathcal{F}}\rightleftharpoons \! \hatbsf$ are a Fourier pair~\eqref{Def:F-pair}. Notice that the summation can be rewritten using a dot product~\eqref{Def:dotProduct}. Then we can use Parseval/Plancherel's theorem~\eqref{Eq:Parseval_fg_dotprod}, the $a$-diminished transform~\eqref{Def:a-diminished_F-pair}, and the (auto)convolution theorem~\eqref{Eq:f^m=*^mfhat} to express this in the Fourier domain. Hence, for integer $m\geq0$ and some integer $0 \leq p \leq m$, the \mth\ general moment of~$\bs{f}$ about point~$a\in\C$ in terms of its Fourier coefficients is
\begin{center}
\fbox{
\addtolength{\linewidth}{-55\fboxsep}
 \begin{minipage}{\linewidth}
    \vspace{-5pt}
    \begin{eqnarray}
    \dimin{\mu}{a}_m\gr{\bs{f}} &:=& \frac{1}{\abs{G}} \sum_{i\in G} \gr{f_i-a}^m \nonumber \\
    &=& \frac{1}{\abs{G}} \Dotprod{\gr{\bs{f} - a\bs{1}}^{m-p}}{\gr{\bs{f} - a\bs{1}}^p} \;\; \nonumber\\[5pt]
    &=& \Dotprod{\!\conv^{m-p} \!\dimin{\hatbsf}{a}}{\conv^p \dimin{\hatbsf}{a}_{\!\ominus}}. \label{Eq:mu_m^a=g*gtg*g}
    \end{eqnarray}
    \vspace{-15pt}
 \end{minipage}
}\\[15pt]
\end{center}

\begin{thm}[$m$-Coefficient, Index Annihilation] \label{Thm:mthMoment_a}
Assume function~$\bs{f}:G\rightarrow\C$ with~$G$ a finite abelian group~\eqref{Eq:FiniteAbelianGroup_G}, and Fourier pair $\bs{f} \stackrel{\mathcal{F}}\rightleftharpoons \hatbsf$~\eqref{Def:F-pair}. Then for integer $m\geq1$, the \mth\ general moment of~$\bs{f}$ about point~$a\in\C$ is of the form
\begin{equation} \label{ThmEq:mthMoment_a}
\dimin{\mu}{a}_m\gr{\bs{f}}
\,\;=\;\; \!\!\sum_{\substack{j_q\:\!\in\:\!G\\[2pt]
\bigoplus_{q=1}^m j_q = 0}} \dimin{\hat{f}}{a}_{j_1} \dimin{\hat{f}}{a}_{j_2} \cdots \dimin{\hat{f}}{a}_{j_m}
\end{equation}
where the $a$-diminished transform~$\dimin{\hatbsf}{a}$
is defined in~\eqref{Def:a-diminished_Xform}.
That is, each term consists of a product of precisely $m$ Fourier coefficients, $\dimin{\hat{f}}{a}_{j_1} \dimin{\hat{f}}{a}_{j_2} \cdots \dimin{\hat{f}}{a}_{j_m}$, such that the modulo summation of the ${j_q}$'s is zero.\footnote{Index $j_q\in G$ identifies the \qth\ element of the Fourier transform. It should not be confused with the \qth\ position of an arbitrary $n$-tuple $j = \vect{\bcode{j}_1,\bcode{j}_2,\ldots,\bcode{j}_n} \in G$, where $\bcode{j}_q\in\Int_{N_q}$.}
\end{thm}

\begin{proof}
Fix some integer $m\geq1$. Examining the right-hand side of~\eqref{Eq:mu_m^a=g*gtg*g} with any integer $0\leq p\leq m$, we see from~\eqref{Eq:C_*^mf=C_f^m} that $\conv^{m-p} \dimin{\hatbsf}{a}$ is just column $j=0$ of~$\bs{C}_{\!\!\dimin{\hatbsf}{a}}^{m-p}$, and $\conv^p \dimin{\hatbsf}{a}_{\!\ominus}$ is row $i=0$ of~$\bs{C}_{\!\!\dimin{\hatbsf}{a}}^{p}$. Thus we have $\dimin{\mu}{a}_m\;\!\!\gr{\bs{f}} = \bs{C}_{\!\!\dimin{\hatbsf}{a}}^m \gr{0,0}$.
Here, the relative shift is $i \ominus j=0$.
It follows that the main diagonal of $\bs{C}_{\!\!\dimin{\hatbsf}{a}}^m$ is
\begin{equation} \label{Eq:C_fC_f_kk-mu_m}
\bs{C}_{\!\!\dimin{\hatbsf}{a}}^m \gr{k,k} \,=\, \dimin{\mu}{a}_m\;\!\!\gr{\bs{f}}
\end{equation}
because for any $i=j=k\in G$, the relative shift will also be $i \ominus j=0$.

At the same time, we can invoke~\eqref{Eq:C_fC_g_ij} recursively to claim for any $k,l\in G$ that the \klth\ element of $\bs{C}_{\!\!\dimin{\hatbsf}{a}}^m = \bs{C}_{\!\!\dimin{\hatbsf}{a}}\bs{C}_{\!\!\dimin{\hatbsf}{a}}\cdots\bs{C}_{\!\!\dimin{\hatbsf}{a}}$ is
$$
\bs{C}_{\!\!\dimin{\hatbsf}{a}}^m \gr{k,l} = \sum_{j_q\in G} \hat{f}_{j_1}^{\gr{a}} \hat{f}_{j_2}^{\gr{a}} \cdots \hat{f}_{j_m}^{\gr{a}} \quad \text{s.t.} \quad j_1 \oplus j_2 \oplus \cdots \oplus j_m = k \ominus l.
$$
But the main diagonal enforces $k=l$, which yields~\eqref{ThmEq:mthMoment_a}.
\end{proof}

\end{document}
